\section{Four campaign case studies}
\label{sec:cases}

\subsection{The decoy-gradient experiment}
\label{sec:f1}

Padding batches with synthetic decoy rows aims to obscure which activation
vectors correspond to real user tokens \cite{paper1}. During each split
step, the trusted client transmits protected forward activations to the
cloud, receives remote hidden states, evaluates the private training loss
locally, and sends output gradients back to the cloud (\artifact{E1}).
Decoy rows receive identically zero gradient. Because the local loss
ignores synthetic positions, that deterministic zero-support mask exposes
the exact real-versus-decoy row partition without revealing token
identities directly, much as gradient-sparsity patterns leak metadata in
federated vision and trajectory settings \cite{gisurvey2022,trace2026}.

Archived logs document exact partition recovery across all 4,096 inspected
frames \emph{per run} over nine training seeds (\artifact{E1}). Recovering
token identities requires a separate content attack. On that distinct token
denominator, a frequent-word probe beats a constant-guess baseline by
$+0.65$ to $+1.50$ percentage points. Across six utility-compliant cells
that expose backward gradients (three seeds on each of two datasets),
forward-activation probes stay inside the $+1.0$ percentage-point privacy
gate fixed prior to execution, whereas joint forward-and-backward
evaluators breach that gate decisively.

Forward-only non-detection remains accurate for an observer restricted to
forward activations. Extending that assurance to the cloud's actual
bidirectional view constitutes an out-of-scope inference that the recorded
backward breach refutes. Clipping each gradient row and adding Gaussian
noise eliminates the scored partition leak across all 1,024 audited frames
(producing $0.400$ partition agreement at a $32/80$ base rate) while
imposing roughly $0.01$ nats of held-out cross-entropy \emph{utility cost}
(\artifact{E1}). That $0.01$-nat figure measures validation loss degradation
on the held-out split, while multi-step gradient accumulation and exact
batch-reconstruction attacks \cite{dager2024,fedspy2026} remain untested
against the noisy channel. Even within that bounded probe scope, closing a
deterministic support leak via per-row clipping and noise refutes blanket
claims that perturbation mechanisms never mitigate training-time leakage.

\subsection{The boundary-replacement audit}
\label{sec:f2}

Auditing the boundary-replacement codebase uncovered an off-by-one indexing
shift between captured representations and target token IDs (\artifact{B1}).
Hidden states for $(x_1,\ldots,x_{T-1})$ met labels $(x_2,\ldots,x_T)$.
Every activation row encoding current token $x_t$ was consequently evaluated
on how accurately a probe predicted the subsequent token $x_{t+1}$. High
canary scores under that shifted alignment validate next-token predictive
signal alone and invalidate claims of resistance to current-token input
reconstruction, though the raw numbers remain interpretable if explicitly
relabelled as next-token probes.

Independent artifact-review records from the same study also re-examined
whether remote computation improved model quality (\artifact{B1}).
Source-reported comparisons across six clean corpus-and-seed pairs show
that standalone client baselines trained from scratch with zero remote
return match or outperform the split-trained remote model every time.
Learned remote layers were consequently unnecessary for those evaluated
endpoints. Zeroing the return pathway of an already trained split
checkpoint tests a different counterfactual, since co-adapted client layers
naturally degrade when their expected remote input disappears mid-inference.

Checkpoint-level scoring exposed an equally sharp divergence between
continuous loss and discrete generation metrics (\artifact{B1}). Continuous
loss cleared its threshold. Held-out negative log-likelihood rose by roughly
$0.034$ nats within rule, yet teacher-forced top-1 match on answer tokens
dropped $16.64$ percentage points from $34.79\%$ to $18.15\%$. Answer-token
matching under teacher forcing differs from standard multi-choice benchmark
accuracy and comes here from a single audited checkpoint. Paired with the
unmeasured current-token reconstruction risk, this discrepancy narrows both
the delegation and utility conclusions supported by the archive.

\subsection{The privacy, utility and gain exploration}
\label{sec:f3}

Across 232 training runs spanning 19 candidate configurations and four
seeds per cell, the exploratory trilemma campaign rejected every
configuration under its frozen rules (\artifact{R1}). Recomputing the cell
summary yields 13 qualified privacy breaches, five utility failures, and
seven remote-gain failures. Rejection causes overlap across cells. All six
configurations whose privacy status remained undetermined failed utility or
learning-gain gates independently, including \texttt{width-64} on gain
alone.

Only four of the 24 required privacy probes passed pre-test qualification.
Those four calibrated instruments consisted of input-byte probes operating
on inbound and outbound gradients across the two evaluated defence families
(\artifact{R1}). Unqualified probes prevent certifying a joint privacy
pass, yet breaches triggered by the four calibrated probes remain
scientifically valid alongside the independent utility and gain failures.

Boundary proximity on the \texttt{span-6} cell illustrates why audit
scripts must check exact floating-point inequalities (\artifact{R1}).
Across four training seeds, \texttt{span-6} recorded a mean utility delta
of $0.0999705344$ nats ($2.95\times 10^{-5}$ nats inside the $0.10$ limit)
and an upper 95\% Student-$t$ bound of $0.1195024415$ nats (inside the
$0.12$ cap). Both utility gates passed. Saved rejection flags instead trace
to learning-gain failure paired with a qualified byte-probe breach across a
four-seed exploratory cohort unadjusted for multiplicity.

An accompanying confirmation proposal specified eight seeds per candidate
(\artifact{R1}). Program leads never froze, authorised, or ran that stage.
Using plug-in exploratory means and variances, a prospective calculation
bounded whole-family confirmation power at $0.263772$ (dropping sharply
under wider variance multipliers) and served strictly as a planning stop
rule for future compute allocation.

\subsection{The defence-redesign pilot}
\label{sec:f4}

Rebuilt probes drove a subsequent seven-cell defence-redesign pilot
comprising 66 runs across three paired seeds per cell and evaluation panels
of 16,384 token windows (\artifact{P4}). Zero candidates achieved pilot
eligibility. Every cell cleared the causal remote-gain gate and five
cleared utility, whereas one-bit and two-bit activation quantisation failed
their utility limits. Under the campaign's all-bindings aggregation policy,
composite privacy assessments recorded four failures and three undetermined
cells.

Calibrated atomic probes nonetheless captured severe leakage across the
tested quantisation and noise mechanisms (\artifact{P4}). Worst-seed
punctuation correlation advantages spanned $0.79$ to $0.89$ against an
acceptance ceiling of $0.02$, while next-byte class accuracy advantages
reached $0.097$ to $0.112$. Unlike Section~\ref{sec:f2}, next-byte class served here as an intentional
target. Observing calibrated probe advantages $40\times$ above a
policy threshold refutes the candidate defences under that rule without
implying probe malfunction. Recent split-LLM obfuscation studies report
comparable structural persistence \cite{privpair2026,mixguard2026}.

Calibration source files separate four numerical constants by role
(\artifact{P4}). Positive-control canary floors sit at $0.10$ for semantic
properties and $0.02$ for structural targets, whereas both the permutation
null band and the candidate acceptance gate equal $0.02$. Passing a canary
floor confirms probe responsiveness on unprotected controls, and
qualification logs record marginal null-band excursions as finite-sample
variance on bounded exploratory panels.
